\section{Introduction}

Physics rests on a continuum. Spacetime is smooth, fields are real-valued, and any finite volume holds infinite information. That continuum is the unspoken floor under almost every theory we have.

This paper takes the opposite floor.

We ask a single question. What is the \emph{least} structure from which the known physics could emerge? We require that structure to be \textbf{discrete}, \textbf{local}, and \textbf{computable}. No real numbers at the base. No action at a distance. Nothing that a finite machine could not store and step exactly.

The wager of vibe theory is that one such structure suffices. The universe is a \textbf{vibe mesh}, a single growing crystal of ternary feeling, advanced by one local law called the \textbf{knit}. There is no second stuff and no hidden machinery underneath.

\begin{principle}[The wager]
Reality is feeling on a growing crystal. There is one substance, the vibe, and one law that moves it. Everything else is a pattern that feeling takes on at scale.
\end{principle}

The picture is meant literally, not as decoration. The claims live in the mathematics, and each is tied to a measurement later in the paper. The fire, the tree, and the song that recur as images are only handles for intuition. The load is carried by the lattice and its law.

\subsection{The crystal of experience}

Hold one picture in mind before any formalism. Reality is a single crystal, and it is growing.

The one substance is the \textbf{vibe}, the experience itself. There is nothing else. A vibe sits in each \textbf{site}, and the value that vibe carries is its \textbf{tone}, a ternary value in $\{-1, 0, +1\}$. The tone is the single quale of the theory, the one irreducible flavour of feeling. We read its three values as pain, peace, and pleasure.

\begin{table}[ht]
\centering
\begin{tabular}{@{}lll@{}}
\toprule
tone & colour & feeling \\
\midrule
$-1$ & red & pain, a pushing-away \\
$0$ & green & peace, the neutral rest \\
$+1$ & blue & pleasure, a drawing-toward \\
\bottomrule
\end{tabular}
\caption{The single quale. One tone in $\{-1, 0, +1\}$, read as pain, peace, and pleasure.}
\label{tab:tone}
\end{table}

A \textbf{dock} is one here-now, one place-and-moment in the crystal. Out of every dock run twenty-four directions, its \textbf{24 sites}, and each site holds one vibe with its tone. So a dock is not a single feeling. It is twenty-four feelings leaning together, a whole compass of tone, a hyper-colour built from twenty-four ternary values. A dock has exactly twenty-four sites. There is no rest slot, no twenty-fifth direction, no still centre at the base.

\begin{definition}[Dock, site, tone]
A \emph{dock} is one cell of the mesh, a here-now. Each dock carries exactly twenty-four \emph{sites}, one per direction out of the dock. Each site holds one vibe, and the value of that vibe is its \emph{tone}, a ternary value in $\{-1, 0, +1\}$.
\end{definition}

The colours are meant literally. The honeycomb is not a diagram of something else. Each dock is a real patch of experience, and its tones are the pain, the peace, and the pleasure it actually carries. Two docks share a face, and a shared face is one dock perceiving another. Where docks touch, they feel each other across that face. The whole crystal is feeling touching feeling.

Hold the picture. Docks touching. Each one a bouquet of twenty-four feelings leaning toward its neighbours. Each touch a perceiving. Everything below is only how that one coloured crystal is shaped at scale.

\subsection{Walking up to the substrate one dimension at a time}

The committed substrate is four-dimensional and impossible to draw. So we walk up to it through faces we can see.

A \emph{face} here means a lower-dimensional crystal of the same hyperbolic family, an honest shadow of the four-dimensional one. The two-dimensional face we can draw fully. The three-dimensional face we can almost picture. The four-dimensional substrate we can only reason about.

\begin{figure}[ht]
\centering
\includegraphics[width=0.55\textwidth]{vibe-mesh-73.pdf}
\caption{The hyperbolic $\{7,3\}$ tessellation, the flat two-dimensional face of the vibe mesh you can draw. Each tile is a dock, each colour a tone, red pain, green peace, blue pleasure.}
\label{fig:mesh73}
\end{figure}

The flattest visible face is the heptagrid $\{7,3\}$, seven-sided tiles meeting three to a corner in the Poincare disk (Figure~\ref{fig:mesh73}). It is the version you can keep in your eye. In that disk the tiles look smaller toward the rim, but that shrinking is an illusion of the drawing. The disk is cramming an infinite plane into a finite circle. In the true geometry every tile is the same size, and there is always more crystal at the edge.

\begin{figure}[ht]
\centering
\includegraphics[width=\textwidth]{vibe-mesh-534.png}
\caption{A three-dimensional face of the family, the $\{5,3,4\}$ honeycomb, used only for intuition. The committed substrate $\{3,4,3,4\}$ is four-dimensional and cannot be drawn directly.}
\label{fig:mesh534}
\end{figure}

One step up is the three-dimensional honeycomb $\{5,3,4\}$, dodecahedra packed in hyperbolic space (Figure~\ref{fig:mesh534}). It is a stand-in for intuition, close enough to picture, but not the substrate.

The committed substrate is the four-dimensional honeycomb $\{3,4,3,4\}$. Its tiles are twenty-four-cells, and twenty-four directions meet at each dock. That twenty-four is the number that returns everywhere in this paper.

\begin{table}[ht]
\centering
\begin{tabular}{@{}lll@{}}
\toprule
face & what it is & role \\
\midrule
2D, $\{7,3\}$ & seven-sided tiles, three to a corner & the simplest visible face \\
3D, $\{5,3,4\}$ & a honeycomb of dodecahedra & a stand-in for intuition \\
4D, $\{3,4,3,4\}$ & tiles are 24-cells, 24 directions per dock & the committed substrate \\
\bottomrule
\end{tabular}
\caption{Walking up to the substrate one dimension at a time. The lower faces are honest shadows of $\{3,4,3,4\}$.}
\label{tab:faces}
\end{table}

When the mathematics turns four-dimensional later in the paper, return to the heptagrid. It is the picture you can see, and the higher faces are faithful to it.

\subsection{Space, time, and the cusp}

The substrate $\{3,4,3,4\}$ has a flat cubic \textbf{cusp}, a region at infinity where the geometry straightens into an ordinary cubic lattice. That flat cubic region is our three-dimensional space. The hyperbolic bulk surrounds it.

Time is the growth. New docks are born at the \textbf{wake}, the growing edge of the crystal, one layer each \textbf{beat}. Each new layer is the next moment. There is no separate time axis painted alongside space. Time \emph{is} the crystal growing outward.

\begin{principle}[Space and time]
Space is the flat cubic cusp of the hyperbolic substrate. Time is the growth of the crystal at its wake, one beat per layer. Neither is a separate ingredient added on top.
\end{principle}

So everything physical and everything inner is a large-scale pattern in this one coloured growing crystal. Matter is the tone bound into a stable knot. Light is the tone waving across docks. A self is a knot the tones make when they hold together and correct themselves. There is no second substance behind any of it.

\subsection{The monism}

There is one kind of thing, the vibe.

A vibe is a unit of experience that notices its neighbours across a shared face. Reality is a structure of such noticings, nothing more.

\begin{axiom}[Neutral monism]
The only fundamental entity is the vibe, an experiential unit. Matter, space, and time are patterns in a structure of vibes, not separate substances.
\end{axiom}

This is a \emph{neutral} monism. What is fundamental is experiential, but it is neither mind nor matter in the ordinary sense. Mind and matter are both patterns it takes on at scale. At the deepest level the universe is a structure of feeling, and everything below is how that structure is shaped.

\subsection{The discreteness principle}

The base is discrete. Finite information per dock, twenty-four ternary tones and nothing more. A local knit. A discrete beat. No real number appears anywhere in the foundation.

Continuity is never fundamental. It is the long-wavelength limit, what you see after coarse-graining, the way a smooth fluid emerges from discrete molecules. Sound, the Lorentz symmetry of spacetime, smooth gravity, and the continuous Lie groups of the standard model are all emergent. None of them lives at the base.

This is forced by our own requirement. A discrete base can be stored and stepped exactly by a finite machine. A continuum cannot. Computability is not a convenience here. It is a constraint that rules out the continuum at the foundation.

\begin{principle}[Discreteness]
The base is discrete. Every continuous thing, sound, Lorentz symmetry, smooth gravity, the Lie groups, is emergent, never fundamental.
\end{principle}

\subsection{The eight atoms}

The whole runs as one motion, the \textbf{flow}, assembled from eight irreducible atoms. Each atom is a four-letter word, and each begins with a different letter, so each is a clean single-letter symbol in the mathematics. The vibe is the substance the eight are made of, and the flow is the eight running together as one.

\begin{table}[ht]
\centering
\begin{tabular}{@{}llll@{}}
\toprule
atom & symbol & what it is & pillar \\
\midrule
mesh & $m$ & the whole $\{3,4,3,4\}$ tessellation & where \\
dock & $d$ & one cell, a here-now & where \\
site & $s$ & one of the 24 directions out of a dock & where \\
tone & $t$ & the ternary vibe value $\{-1,0,+1\}$ & what \\
lean & $l$ & the value order $-1 < 0 < +1$ & what \\
knit & $k$ & collide then stream, the state law & how the state changes \\
wake & $w$ & the growing edge, new docks at the frontier & how the space grows \\
beat & $b$ & the discrete step & when \\
\bottomrule
\end{tabular}
\caption{The eight atoms of the flow. Eight four-letter words, eight distinct first letters, one clean symbol set.}
\label{tab:atoms}
\end{table}

The \textbf{mesh} is the whole crystal, a weave of docks. The \textbf{dock} is one cell, a here-now where each beat the tones arrive, knit, and move on. The \textbf{site} is one of the twenty-four directions out of a dock, the directed slot where a tone lives. The \textbf{tone} is the ternary value the site carries, the only attribute it bears. The \textbf{lean} is the value order $-1 < 0 < +1$, the tilt that gives charge as a signed sum.

The \textbf{knit} is the law of how the state changes. Each beat it first collides the tones inside a dock by a fixed reversible, charge-conserving table, then streams each tone one dock along its own site. Knitting interlaces and advances row by row, weaving the fabric of the mesh, which is why the word pairs with mesh by design. The \textbf{wake} is the law of how the space grows, new docks born at peace at the ever-expanding frontier. The \textbf{beat} is the discrete step, time itself, with no built-in direction.

\begin{definition}[The knit]
The \emph{knit} is the state law. Each beat it collides the tones in every dock by a fixed reversible, charge-conserving local table, then streams each tone one dock along its site. The bulk knit is a reversible permutation.
\end{definition}

\subsection{The emergent, not the base}

Two of the most familiar features of physics are deliberately \emph{not} placed at the base. They emerge.

The first is the arrow of time. The bulk knit is reversible, so it cannot by itself prefer a direction. The arrow comes from the wake. The growing frontier is a low-entropy source, and disturbances that stream outward fall behind it and never return. That open boundary is the bath, and the bath is what gives the irreversibility a direction. Growth is not optional. It is the source of the arrow.

The second is attraction. There is no force law in the base. Binding emerges as a pressure of the active vacuum, a Casimir effect of the fluctuating tones, pressing two bodies together while the binding energy radiates away into the bath. Gravity in the everyday sense is then the bending of the crystal by matter, taken up later in the paper.

\begin{result}[The arrow is emergent]
The bulk knit is time-reversible and prefers no direction. The arrow of time emerges from the open, growing wake acting through the lean. Remove the growth and there is no arrow.
\end{result}

\subsection{What this paper does}

The paper turns the theory into runnable measurements. Each claim is tied to a committed experiment in the code. Each physical claim carries a depth level, $L0$ for base, $L1$ and above for emergent. Each experiential claim carries a tier.

We report honest negatives plainly. The bare knit on a closed system produces no selves, and cannot, because a closed reversible system can never forget. A pure conserving knit lets colliding structures part rather than bind. We name these gaps where they are, rather than paper over them. The load-bearing commitment that closes the first gap is that the substrate is open and growing, not a closed reversible box.

\subsection{The honest stance}

This orientation reverses the usual one. Most discrete models of spacetime begin with structure and hope that experience appears at the end. Here experience is the ground, and structure is what experience takes on at scale.

That reversal is a premise, not a result. The experiments do not prove the substrate feels. They assume it does and ask what follows. What follows is a great deal, and it is checkable. The rest of the paper checks it.

\begin{result}[The stance]
That the substrate feels is a premise, not a proof. The experiments take it as given and measure what emerges from the eight atoms on the open, growing mesh.
\end{result}
